% Agent-written proof note, not yet Jorn-reviewed.
%
% Purpose: record the bounded-vertex interpretation of near-singular 4-facet
% systems seen by experiments/dev-quadratic-program.
%
% Labels defined here:
%   def:bounded-relevant-vertices,
%   lem:singular-system-bounded-exclusion,
%   lem:perturbed-bounded-exclusion,
%   prop:near-singular-no-bounded-vertex-certificate,
%   rem:f64-near-singular-experiment
%
% The linear algebra statements are elementary. The bridge to f64 diagnostics
% still needs an outward-rounded residual/minimum-norm implementation before it
% becomes a machine-checkable certificate.

\providecommand{\R}{\mathbb{R}}
\providecommand{\one}{\mathbf{1}}
\providecommand{\dist}{\operatorname{dist}}
\providecommand{\im}{\operatorname{im}}

\subsection{Near-singular four-facet systems under a bounded-vertex assumption}
\label{sec:f64-near-singular-vertices}

Let
\[
  K(A)=\{x\in\R^4:\langle a_i,x\rangle\le 1,\ i=1,\ldots,F\},
\]
where the rows $a_i\in\R^4$ are the dual vertices used by
\texttt{experiments/dev-quadratic-program}.  For a four-subset
$I=\{i_1,i_2,i_3,i_4\}$, write $A_I$ for the $4\times4$ matrix with rows
$a_{i_j}$.

The f64 experiment treats a small determinant of $A_I$ as a near-singular
vertex event.  The determinant alone is not enough information: a singular
system can be inconsistent, can have only distant solutions, or can have
bounded solutions.  The relevant local question is whether the active equations
\[
  A_Ix=\one
\]
can contain a bounded vertex candidate.

\begin{unverified}
\begin{definition}[Bounded relevant vertices]\label{def:bounded-relevant-vertices}
Fix $B>0$.  We say that $K(A)$ satisfies the $B$-bounded relevant vertex
assumption if every vertex of $K(A)$ that may be used by the capacity
combinatorics has Euclidean norm at most $B$.

Equivalently for the local certificate below, a four-facet tuple $I$ may be
discarded from bounded vertex enumeration only when
\[
  \{x\in\R^4:A_Ix=\one,\ \|x\|\le B\}=\varnothing .
\]
The global assumption is the statement that no capacity-relevant vertex is lost
by imposing this bounded search region.
\end{definition}
\end{unverified}

\begin{unverified}
\begin{lemma}[Exact singular-system bounded exclusion]
\label{lem:singular-system-bounded-exclusion}
Let $A\in\R^{4\times4}$ and let
\[
  r(A)=\dist(\one,\im A)
       =\min_{x\in\R^4}\|Ax-\one\|.
\]
If $r(A)>0$, then $Ax=\one$ has no solution.  If $r(A)=0$ and
\[
  m(A)=\min\{\|x\|:Ax=\one\}
\]
satisfies $m(A)>B$, then $Ax=\one$ has no solution with $\|x\|\le B$.
\end{lemma}
\end{unverified}

\begin{proof}
The inequality $r(A)>0$ is exactly the statement that $\one$ is not in the image
of $A$, hence no solution exists.  If $r(A)=0$, the solution set is nonempty,
closed, and convex.  The number $m(A)$ is the distance from the origin to this
solution set.  If $m(A)>B$, every solution has norm larger than $B$.
\end{proof}

\begin{unverified}
\begin{lemma}[Perturbed bounded exclusion]\label{lem:perturbed-bounded-exclusion}
Let $A,E\in\R^{4\times4}$, and suppose $\|E\|\le\varepsilon$ in operator norm.
If there is $x$ with
\[
  (A+E)x=\one,\qquad \|x\|\le B,
\]
then
\[
  r(A)=\dist(\one,\im A)\le \varepsilon B.
\]
Consequently, if $r(A)>\varepsilon B$, no perturbation $A+E$ with
$\|E\|\le\varepsilon$ has a solution of norm at most $B$.
\end{lemma}
\end{unverified}

\begin{proof}
From $(A+E)x=\one$ we get
\[
  \|Ax-\one\|=\|Ex\|\le \|E\|\,\|x\|\le \varepsilon B.
\]
Since $r(A)$ is the minimum of $\|Ay-\one\|$ over all $y$, it is at most
$\|Ax-\one\|$.
\end{proof}

\begin{unverified}
\begin{proposition}[No near-singular indeterminacy under bounded exclusion]
\label{prop:near-singular-no-bounded-vertex-certificate}
Fix $B>0$.  Let $I$ be a four-facet tuple and let $A=A_I$.
Assume the $B$-bounded relevant vertex assumption from
Definition~\ref{def:bounded-relevant-vertices}.  If either
\[
  r(A)>0
\]
or
\[
  r(A)=0\quad\text{and}\quad m(A)>B,
\]
then the tuple $I$ cannot contribute a capacity-relevant bounded vertex.
Therefore an f64 implementation may resolve this near-singular tuple as
\textsc{no bounded vertex} once it has a certified lower bound proving one of
the displayed alternatives.

Moreover, if the input matrix is allowed to move by operator norm at most
$\varepsilon$ and $r(A)>\varepsilon B$, then no such perturbation can create a
bounded solution for the tuple $I$.
\end{proposition}
\end{unverified}

\begin{proof}
The exact-input statement is Lemma~\ref{lem:singular-system-bounded-exclusion}
applied to $A_Ix=\one$.  Under the bounded relevant vertex assumption, any
capacity-relevant vertex has norm at most $B$.  If no solution with norm at most
$B$ exists, discarding this four-tuple cannot remove a capacity-relevant vertex.

The perturbative statement is Lemma~\ref{lem:perturbed-bounded-exclusion}.
\end{proof}

\begin{remark}[Connection to the f64 capacity experiment]
\label{rem:f64-near-singular-experiment}
The experiment command
\[
  \texttt{cargo run -p exp-dev-quadratic-program --release --bin f64-capacity-near-singular}
\]
computes diagnostic versions of $r(A_I)$ and the least-norm least-squares
solution for near-singular four-tuples.  These diagnostics are not yet formal
certificates because their residuals and norms are computed in ordinary f64
without outward rounding.  They are evidence for which branch of
Proposition~\ref{prop:near-singular-no-bounded-vertex-certificate} should be
certified.

For the current smoke data with $B=10^3$, the near-singular HKO tuples have
least-squares residuals on the order of $10^{-1}$ to $1$, so
Lemma~\ref{lem:perturbed-bounded-exclusion} says that creating a bounded
solution would require perturbations of the active $4\times4$ matrix on the
order of at least $10^{-4}$ to $10^{-3}$ in operator norm.  The closest
ascent-product examples have much smaller residual but least-norm solution
scale about $10^{12}$, so they fall outside the $B=10^3$ bounded region.
\end{remark}
